\subsection{R-33: generic projected permanents have finite projective stabilizer}

\paragraph{Statement R-33.}
For $m\geq17$ and $4\leq q<m^2$, a general form in the
projected permanent cone $\mathcal P_{m,q}$ has finite complex
projective stabilizer. Consequently, for any positive integer
$s$ and any auxiliary Cartan weight $\theta$, the enhanced
cone $X$ built from this image has
\begin{equation}\label{R33:quotient_upper}
 \dim\mathbb C[X]^{\mathrm{SL}_q}
 \leq \dim X-(q^2-1)
 =q m^2-2m-q^2+3+r_\theta.
\end{equation}
The conclusion is a quotient-dimension upper bound. It makes
no assertion that a generic orbit is closed, that a stable
point exists, or that the nonzero invariant degrees have gcd
one.

\begin{proof}
Write $N=m^2$, and consider the coefficient map
\[
 \Phi:(\operatorname{Mat}_m)^q\longrightarrow
               \operatorname{Sym}^m\mathbb C^q,
 \qquad
 (A_1,\ldots,A_q)\longmapsto
             f_A(x)=\operatorname{per}\!\left(\sum_i x_iA_i\right).
\]
By G-21, its generic differential rank is $qN-2m+2$.
Its generic kernel therefore has dimension $2m-2$.
The row-column diagonal variations
\begin{equation}\label{R33:gauge}
 \delta A_i=\operatorname{diag}(u)A_i+A_i\operatorname{diag}(v),
                   \qquad \sum_a u_a+\sum_b v_b=0,
\end{equation}
always lie in that kernel. On the nonempty open set where
one coefficient matrix has no zero entry, these variations
have dimension exactly $2m-2$: the only redundancy is
$u_a=t$, $v_b=-t$ for every $a,b$. Thus on the intersection
with the maximal-rank open they give the entire kernel.

Independently require $A_1,\ldots,A_q$ to be linearly
independent and their span $W\subset\operatorname{Mat}_m$
to have every Pl\"ucker coordinate nonzero in the matrix-unit
basis. This is another nonempty open: all-coordinate Pl\"ucker
nonvanishing is a finite intersection of nonempty minor opens
on the irreducible Grassmannian, and independent frames map
surjectively to that Grassmannian. The indicated parameter
opens therefore have nonempty intersection. Fix a pencil
in this intersection with $f_A\ne0$.

Let $H\in\mathfrak{gl}_q$ be an infinitesimal projective
symmetry, so $Hf_A=c f_A$. Subtract $(c/m)I_q$ from $H$;
the resulting $H_0$ acts trivially on $f_A$, since scalar
variable multiplication acts on a degree-$m$ form by $m$.
The corresponding coefficient variation is a linear
combination of the $A_i$ and lies in $\ker D\Phi_A$.
By \eqref{R33:gauge}, there are $u,v$ with the stated sum
condition such that this variation is obtained by applying
the diagonal operator
\[
 D(E_{ab})=(u_a+v_b)E_{ab}
\]
to every $A_i$. In particular $D(W)\subseteq W$.

The induced operator on the line $\bigwedge^q W$ is scalar.
Since every Pl\"ucker coordinate is nonzero, the sum of the
$D$-weights over every $q$-element subset of the $N$ matrix
units is the same. Because $1\leq q<N$, any two matrix-unit
indices can be exchanged while retaining $q-1$ other indices.
Their individual weights must therefore be equal. Thus
$u_a+v_b=\tau$ for every $a,b$. Summing all $N$ weights gives
\[
 N\tau=m\left(\sum_a u_a+\sum_b v_b\right)=0,
\]
so $D=0$. The $A_i$ are independent, hence the coefficient
variation vanishes only when $H_0=0$. We have proved that
every infinitesimal projective symmetry is scalar.

The projective algebraic stabilizer of $[f_A]$ therefore has
zero Lie algebra and is finite in characteristic zero.
Finiteness is open here, since the Lie stabilizer is the
kernel of a matrix whose entries are polynomial in the form
coefficients. The nonempty open just established meets the
image, so the general point of $\mathcal P_{m,q}$ has this
property.

For the enhanced cone, a projective symmetry of a pure tensor
$F^{\otimes s}\otimes v_\theta$ preserves $[F]$. Its stabilizer
inside $\mathrm{SL}_q$ is therefore finite over the generic
finite-projective-stabilizer locus just proved. Thus the
generic orbit in $X$ has dimension $q^2-1$.
The quotient morphism
\[
 X\longrightarrow\operatorname{Spec}\mathbb C[X]^{\mathrm{SL}_q}
\]
is dominant. Its generic fiber meets the dense finite-stabilizer
open and contains the orbit through such a point. Hence the
generic fiber has dimension at least $q^2-1$.
The generic fiber dimension theorem gives
\eqref{R33:quotient_upper}. Here
$\dim X=\dim\mathcal P_{m,q}+r_\theta$ by the same
Segre--Veronese--Cartan construction as G-22, and G-21 gives
$\dim\mathcal P_{m,q}=q m^2-2m+2$.
\end{proof}

\paragraph{Sources and scope.}
The new argument uses the proved exact dimension and generic
kernel statement from G-21, then elementary Pl\"ucker weights.
The final generic-fiber equality is the Stacks Project,
Lemma~33.20.4(2), tag~\texttt{0B2L} (R-B18), whose stated
proof uses transcendence degrees. The EGA/SGA~3 framework
in dictionary entry D-S3 concerns fiber-dimension upper
semicontinuity and the stabilizer-openness framework; it
is not substituted for this generic-fiber equality.
No equality with the invariant quotient dimension is inferred
from generic freeness. A ring used only for an upper Hilbert
bound can use this inequality even when its stable locus or
scalar period is unknown. The restriction $q<m^2$ is essential
to the all-subset Pl\"ucker exchange argument and is retained.
